\begin{textsample}{POS Dim 2 – 2021\_c4 – Score 20.00 – 2021\_c4/enade\_2021\_letras\_por\_d96b21bbf0.txt}  \label{ex:f2_pos_034}
CAPÍTULO 4 PERCEPÇÃO DA PROVA Nas análises feitas neste capítulo, constam as percepções dos concluintes da Área de Letras-Português ( Licenciatura ) sobre a prova aplicada no Enade/2021.
Essas percepções foram mensuradas por meio de nove questões que avaliaram desde o grau de dificuldade da prova até o tempo gasto para concluí -la.
As percepções sobre a prova foram relacionadas ao desempenho dos estudantes, à Grande Região de funcionamento do curso, à Categoria Administrativa e ao tipo de Organização Acadêmica da IES.
O \textbf{questionário} de percepção da prova encontra -se ao final do Anexo VIII, onde está a \textbf{reprodução} do exame.
O desempenho dos estudantes foi classificado em quatro quartos.
Para tanto, esse desempenho foi \textbf{ordenado} de forma ascendente.
O percentil 25, P25, também conhecido como primeiro quartil, é a nota de desempenho que deixa um quarto ( 25 \% ) dos valores observados abaixo e três quartos acima.
Na Figura 1, apresenta -se uma ilustração desse conceito.
O quarto inferior de desempenho é composto pelas notas abaixo do primeiro quartil.
Já o percentil 75, P75, também conhecido como terceiro quartil, é o valor para o qual há três quartos ( 75 \% ) dos dados abaixo e um quarto acima dele.
O quarto superior de desempenho é composto pelas notas iguais ou acima do terceiro quartil.
O percentil 50, P50, também conhecido como mediana, é o valor que divide as notas em dois conjuntos de igual tamanho.
O segundo quarto inclui os valores entre o primeiro quartil ( P25 ) e a mediana.
O terceiro quarto contém os valores entre a mediana ( P50 ) e o terceiro quartil ( P75 ).
Vale ressaltar que percentis, quartis e medianas são pontos que, não obrigatoriamente, pertencem ao conjunto original de dados, ao passo que os quartos são \textbf{subconjuntos} dos dados originais.
A seguir, serão apresentados gráficos com resultados selecionados, relativos às nove questões avaliadas por grupos de estudantes.
Nas barras dos gráficos, apresenta -se o percentual de alunos que assinalaram uma das opções ou a soma das porcentagens daqueles que assinalaram duas ( ou três ) delas.
Por exemplo, para as questões 1 e 2, nos gráficos, é apresentada a porcentagem total de participantes que assinalaram as opções ( D ) Difícil e ( E ) Muito difícil.
Em cada barra, foram assinalados, também, os extremos do intervalo de confiança de 95 \% como linhas verticais unidas por uma linha horizontal na forma da letra H maiúscula.
O \textbf{estimador} de um \textbf{parâmetro} com um certo nível de confiança ( por exemplo, 95,0 \% ) deve conter o \textbf{parâmetro} no intervalo de confiança em 95 \% das vezes.
Na comparação entre os \textbf{estimadores} dos \textbf{parâmetros} de duas classes de uma dada categoria ( por exemplo : Norte e Nordeste nas Grandes Regiões, de primeiro e último quarto dentro de desempenho, de IES Públicas e Privada, ou de Universidades e Faculdades ), associados aos seus respectivos intervalos de confiança, diz -se que não há diferença estatisticamente significativa entre os \textbf{parâmetros} das duas categorias quando há \textbf{interseção} entre os intervalos de confiança, e que há diferença, se os intervalos de confiança são \textbf{disjuntos} ( para maiores informações \textbf{vide} Glossário ).
Nas Tabelas do Anexo II, são apresentados os valores absolutos e a distribuição percentual18 das alternativas válidas das nove questões, segundo o mesmo \textbf{recorte} de desempenho dos alunos e da Grande Região de funcionamento do curso, de Categoria Administrativa e do tipo de Organização Acadêmica da IES. 4.1 GRAU DE DIFICULDADE DA PROVA 4.1.1 COMPONENTE DE FORMAÇÃO GERAL Ao avaliarem “ Qual o grau de dificuldade desta prova na parte de Formação Geral? ” ( Questão 1 ), 35,1 \% do grupo de inscritos e presentes optaram pelas alternativas Difícil ou Muito difícil.
Entretanto, para mais da metade dos estudantes ( 54,8 \% ), o Componente de Formação Geral da prova foi avaliado com grau de dificuldade Médio ( Gráfico 4.1, Gráfico 4.2 e, no Anexo II, a Tabela II.1 ).
O percentual de estudantes que consideraram a prova Difícil ou Muito difícil foi maior na região Sul, onde a proporção foi de 41,4 \%, enquanto a de menor proporção foi na Sudeste, com 30,0 \%.
No Gráfico 4.1, é possível observar que as diferenças entre a região Sudeste e as regiões Norte, Nordeste e Sul são estatisticamente significativas, assim como a diferença entre a região Sul e as regiões Centro-Oeste e Nordeste.
Nas Grandes Regiões, a proporção de presentes à prova que avaliaram o Componente de Formação Geral como de grau de dificuldade Médio esteve entre 52,0 \%, na região Sul, e 56,1 \%, na região Nordeste. 18 Cumpre lembrar uma das convenções para tabelas \textbf{numéricas} ( ao final do \textbf{Sumário} ) sobre a possibilidade de a soma das partes não resultar em 100 \% por questões de \textbf{arredondamento}. proporção de presentes à prova que avaliaram o Componente de Formação Geral como de grau de dificuldade Médio esteve entre 52,0 \%, na região Sul, e 56,1 \%, na região Nordeste.
O percentual de alunos que consideraram o componente de Formação Geral da prova Difícil ou Muito difícil foi decrescente em função dos quartos de desempenho : 45,3 \%, no primeiro quarto, e 20,9 \%, no quarto, o grupo de melhor desempenho na prova.
Nos quartos de desempenho intermediários, a proporção de alunos que consideraram a prova Difícil ou Muito difícil foi de 41,6 \%, no segundo quarto, e de 33,2 \%, no terceiro quarto.
As diferenças entre todos os quartos de desempenho são estatisticamente significativas.
Para todos os quartos de desempenho, a alternativa modal para essa pergunta foi a Médio, com 46,3 \% e 62,4 \% dos respondentes nos quartos extremos, primeiro e quarto, respectivamente.
O percentual de estudantes que consideraram o componente de Formação Geral da prova Difícil ou Muito difícil foi maior entre os estudantes de Centros Universitários, alcançando 43,1 \%, enquanto entre os estudantes das Universidades, a Organização Acadêmica com a menor incidência, o percentual alcançou 32,8 \%.
No Gráfico 4.3, é possível observar que as diferenças entre Centros Universitários e as demais Organizações Acadêmicas são estatisticamente significativas.
Nas Organizações Acadêmicas, a proporção de presentes à prova que avaliaram o Componente de Formação Geral como de grau de dificuldade Médio esteve entre 50,1 \%, nos Centros Universitários, e 57,5 \%, nas Faculdades ( ver também a Tabela II.2, no Anexo II ).
O percentual de estudantes que consideraram o componente de Formação Geral da prova Difícil ou Muito difícil foi maior entre os estudantes de cursos de IES Privadas ( 41,4 \% ), com diferença estatisticamente significativa da proporção registrada em IES Públicas ( 31,8 \% ).
Já a proporção de presentes à prova que avaliaram o Componente de Formação Geral como de grau de dificuldade Médio foi de 56,5 \%, entre os estudantes de IES Públicas, e de 51,6 \%, entre os de IES Privadas ( ver Gráfico 4.4 ). 4.1.2 COMPONENTE DE CONHECIMENTO ESPECÍFICO Ao responderem à Questão 2 – “ Qual o grau de dificuldade desta prova na parte de Componente Específico? ” – 40,5 \% do grupo de estudantes classificaram -na como Difícil ou Muito difícil.
Além disso, o Componente modal de Conhecimento Específico da prova foi considerado de grau de dificuldade Médio por 51,9 \% dos alunos ( Gráfico 4.5, Gráfico 4.6, e, no Anexo II, a Tabela II.3 ).
A análise das respostas dos estudantes quanto ao grau de dificuldade do Componente de Conhecimento Específico da prova, agregado por Grande Região, mostra diferenças estatisticamente significativas entre proporções de alunos que a avaliaram como Difícil ou Muito difícil.
Tais diferenças são observadas entre a região com a maior proporção, a Sul ( 48,0 \% ), e todas as demais regiões : Norte ( 40,3 \% ), Nordeste ( 38,0 \% ), Sudeste ( 36,5 \% ) e Centro-Oeste ( 37,5 \% ).
O percentual de alunos que classificaram o grau de dificuldade como Médio, no Componente de Conhecimento Específico, variou de 46,9 \% a 54,6 \%, nas regiões Sul e Nordeste, respectivamente.
Considerando -se a avaliação da dificuldade das questões do Componente de Conhecimento Específico da prova, de acordo com o desempenho dos estudantes, observa -se uma tendência decrescente dos que consideraram a parte específica Difícil ou Muito difícil, variando de 30,5 \% ( último quarto ) a 47,7 \% ( primeiro quarto ).
As demais proporções foram de 43,5 \% e 40,6 \%, respectivamente, no segundo e terceiro quartos.
Observa -se que há diferença estatisticamente significativa dos resultados entre o quarto superior de desempenho e os demais, assim como entre o quarto inferior e os demais quartos.
Já os que responderam que o grau de dificuldade das questões do Componente de Conhecimento Específico da prova foi Médio atingiram \textbf{percentagens} que variaram de 44,7 \%, no primeiro quarto, a 59,0 \%, no último.
Considerando -se a avaliação da dificuldade das questões do Componente de Conhecimento Específico da prova, o percentual de \textbf{estudantes} que as consideraram Difícil ou Muito difícil foi maior entre os estudantes de Centros Universitários, atingindo 50,4 \% ; já entre os de CEFET/IF, a Organização Acadêmica com a menor incidência, essa proporção foi de 37,3 \%.
No Gráfico 4.7, é possível observar que a diferença entre os Centros Universitários e todas as demais Organizações Acadêmicas é estatisticamente significativa.
Nas Organizações Acadêmicas, a proporção de presentes à prova que consideraram o Componente de Conhecimento Específico como de grau de dificuldade Médio esteve entre 44,8 \%, nos Centros Universitários, e 54,0 \%, nas Universidades ( ver também Tabela II.4, no Anexo II ).
O percentual de estudantes que consideraram o Componente Específico da prova Difícil ou Muito difícil foi maior entre os estudantes de cursos de IES Privadas ( 47,8 \% ), com diferença estatisticamente significativa da proporção registrada em IES Públicas ( 36,7 \% ).
Já a proporção de presentes à prova que avaliaram esse componente como de grau de dificuldade Médio foi de 54,6 \%, entre os estudantes de IES Públicas, e de 46,8 \% entre os de IES Privadas ( ver Gráfico 4.8 e Tabela II.4, no Anexo II ). 4.2 EXTENSÃO DA PROVA EM RELAÇÃO AO TEMPO TOTAL \textbf{Indagados} quanto à extensão da prova, em relação ao tempo total oferecido para a sua resolução ( Questão 3 ), os estudantes apontaram, com maior incidência, a alternativa que considerava a extensão Adequada, em todas as agregações consideradas ( Gráfico 4.9, Gráfico 4.10, e, no Anexo II, a Tabela II.5 ).
O percentual de alunos que responderam ser a extensão da prova Adequada foi de 38,5 \%.
Já 51,4 \% dos inscritos presentes consideraram que a prova foi Longa ou Muito longa, e 10,1 \% a avaliaram como Curta ou Muito curta.
Entre as Grandes Regiões, os que avaliaram a prova como Longa ou Muito longa, em relação ao tempo total destinado à sua resolução, atingiram uma porcentagem que variou de 41,9 \%, na região Norte, até 53,9 \%, na região Sudeste.
Há diferença estatisticamente significativa entre a região Norte e as regiões Nordeste, Sudeste e Sul.
Considerando -se o desempenho dos alunos, nota -se ainda que 40,5 \% consideraram a extensão da prova Adequada, no quarto de desempenho inferior, e 39,5 \%, no de melhor desempenho ( quarto superior ).
Nos quartos intermediários, essa proporção foi de 37,0 \%, no segundo quarto, e de 36,8 \%, no terceiro.
No Gráfico 4.10, observa -se que as proporções de estudantes que consideraram a prova Longa ou Muito longa apresentaram tendência crescente em função dos quartos de desempenho, sendo a maior proporção dessa resposta no terceiro quarto de desempenho ( 53,6 \% ), e a menor, no primeiro quarto inferior ( 47,7 \% ).
Pode -se constatar que há diferença estatisticamente significativa entre a proporção do quarto inferior de desempenho e a dos demais.
Considerando -se a extensão da prova, o percentual de estudantes que avaliaram a prova como Longa ou Muito longa, em relação ao tempo total destinado à sua resolução, foi maior entre os estudantes de CEFET/IF, alcançando 58,3 \%, enquanto entre os estudantes de Faculdades, a Organização Acadêmica com a menor incidência, esse índice atingiu 47,4 \% ; já entre os de Centros Universitários, o índice alcançou 50,1 \%, e entre os de Universidades, 51,3 \%.
No Gráfico 4.11, é possível observar que há diferença estatisticamente significativa entre os CEFET/IF e os três demais tipos de Organização Acadêmica.
A proporção de presentes à prova que consideraram a sua extensão Adequada, em relação ao tempo total destinado à sua resolução, esteve entre 30,3 \%, nos CEFET/IF, e 41,3 \%, nos Centros Universitários ( ver também Tabela II.6, no Anexo II ).
O percentual de estudantes que consideraram a prova Longa ou Muito longa, em relação ao tempo total destinado à sua resolução, foi maior entre os estudantes de cursos de IES Públicas ( 52,1 \% ), não havendo diferença estatisticamente significativa da proporção em relação aos de IES Privadas ( 50,2 \% ).
Já a proporção de presentes à prova que consideraram a extensão Adequada foi de 37,4 \%, entre os estudantes de IES Públicas, e de 40,5 \%, entre os de IES Privadas ( ver Gráfico 4.12 e Tabela II.6, no Anexo II ). 4.3 COMPREENSÃO DOS ENUNCIADOS DAS QUESTÕES 4.3.1 COMPONENTE DE FORMAÇÃO GERAL Em relação aos enunciados das questões do Componente de Formação Geral ( Questão 4 ), as opiniões foram positivas, já que 76,1 \% dos alunos avaliados consideraram todos ou a maioria dos enunciados das questões “ claros e objetivos ” ( Gráfico 4.13, Gráfico 4.14 e, no Anexo II, a Tabela II.7 ).
Na análise regional, a \textbf{percentagem} de estudantes que avaliaram que todos ou a maioria dos enunciados das questões do Componente de Formação Geral estavam “ claros e objetivos ” variou de 74,5 \%, na região Norte, a 78,3 \%, na região Centro-Oeste.
A análise das percepções dos estudantes sobre a clareza e a objetividade dos enunciados permite afirmar que todos ou a maioria dos enunciados de questões relativas ao Componente de Formação Geral, foram considerados “ claros e objetivos ” para a maior parte dos respondentes ( maior ou igual a 74,5 \%, em todas as regiões, e maior ou igual a 64,2 \%, em todos os quartos de desempenho ).
Segundo o desempenho, observa -se que a proporção dos que emitiram essa opinião \textbf{cresce}, conforme o desempenho aumenta, com diferenças estatisticamente significativas entre todos os quartos de desempenho.
No quarto superior, a clareza e a objetividade de todos ou da maioria dos enunciados das questões foram percebidas por 87,2 \%.
A proporção dos presentes que consideraram que todos ou a maioria dos enunciados das questões do Componente de Formação Geral estavam “ claros e objetivos ” foi maior entre os estudantes de Universidades, chegando a 77,5 \%, enquanto entre os de Faculdades, a Organização Acadêmica com a menor incidência, esse índice atingiu 69,7 \% ; já entre os de Centros Universitários, o índice atingiu 73,0 \%, e entre os de CEFET/IF, 73,1 \%.
No Gráfico 4.15, é possível observar que as diferenças entre as Universidades e os três demais tipos de Organização Acadêmica são estatisticamente significativas ( ver também Tabela II.8, no Anexo II ).
O percentual de estudantes que consideraram que todos ou a maioria dos enunciados das questões do Componente de Formação Geral estavam “ claros e objetivos ” foi maior entre os estudantes de cursos de IES Públicas ( 77,1 \% ), com diferença estatisticamente significativa da proporção registrada em IES Privadas ( 74,4 \% ).
Um maior detalhamento pode ser encontrado no Gráfico 4.16 e na Tabela II.8, no Anexo II. 4.3.2 COMPONENTE DE CONHECIMENTO ESPECÍFICO Em relação aos enunciados das questões do Componente de Conhecimento Específico da prova, para 75,9 \% dos estudantes avaliados da Área de Letras-Português ( Licenciatura ), a clareza e a objetividade ( Questão 5 ) estavam presentes em todas ou na maioria das questões ( Gráfico 4.17 a Gráfico 4.20, e no Anexo II, as Tabelas II.9 e II.10 ).
A maioria dos estudantes de todas as Grandes Regiões brasileiras considerou “ claros e objetivos ” todos ou a maioria dos enunciados das questões do Componente de Conhecimento Específico da prova, percentual sempre maior ou igual a 74,1 \%.
A diferença entre o menor percentual, o da região Nordeste ( 74,1 \% ) e o percentual da região Sudeste ( 77,8 \% ) é estatisticamente significativa.
A proporção de estudantes que consideraram os enunciados das questões “ claros e objetivos ” apresenta uma tendência crescente em relação ao aumento de desempenho : mais elevada no quarto superior ( 86,9 \% ), se comparada ao quarto inferior de desempenho ( 63,9 \% ).
As diferenças entre todos os quartos de desempenho são estatisticamente significativas.
A proporção dos presentes que consideraram que todos ou a maioria dos enunciados das questões do Componente de Conhecimento Específico estavam “ claros e objetivos ” foi maior entre os estudantes de Universidades, atingindo 77,0 \%.
Entre os estudantes de Faculdades, a Organização Acadêmica com a menor incidência, essa proporção atingiu 71,2 \%, enquanto entre os de CEFET/IF atingiu 71,7 \%, e entre os de Centros Universitários, 73,8 \%.
No Gráfico 4.19, é possível observar que as diferenças entre as Universidades em relação a Centros Universitários e CEFET/IF são estatisticamente significativas ( ver também Tabela II.10, no Anexo II ).
O percentual de estudantes que consideraram que todos ou a maioria dos enunciados das questões do Componente de Conhecimento Específico estavam “ claros e objetivos ” foi maior entre os estudantes de cursos de IES Públicas ( 76,5 \% ), sem diferença estatisticamente significativa da proporção alcançada em IES Privadas ( 74,8 \% ).
Ver Gráfico 4.20 e Tabela II.10, no Anexo II, para maior detalhamento. 4.4 SUFICIÊNCIA DAS INFORMAÇÕES/INSTRUÇÕES FORNECIDAS Ao se avaliarem as informações/instruções fornecidas para a resolução das questões ( Questão 6 ), constatou -se que 83,5 \% dos respondentes da Área de Letras-Português ( Licenciatura ) de todo o Brasil afirmaram que eram até excessivas ou suficientes em todas ou na maioria das questões ( Gráfico 4.21, Gráfico 4.22 e, no Anexo II, a Tabela II.11 ).
Quanto à distribuição de respondentes pelas Grandes Regiões, observa -se que a proporção de estudantes que consideraram as informações/instruções fornecidas até excessivas ou suficientes em todas ou na maioria das questões foi sempre superior ou igual a 79,6 \%, chegando a 86,4 \%, na região Sudeste.
Existem diferenças estatisticamente significativas entre a região Sudeste e as regiões Norte, Nordeste e Sul.
Levando -se em conta o desempenho dos participantes, observa -se uma tendência crescente das proporções de participantes que avaliaram as informações/instruções como até excessivas ou suficientes em todas ou na maioria das questões, havendo diferenças estatisticamente significativas entre todos os quartos de desempenho, conforme mostrado no Gráfico 4.22.
O percentual foi mais elevado no quarto superior ( 92,3 \% ), percentual superior à média nacional ( 83,5 \% ).
Já no quarto inferior, a suficiência das informações/instruções declaradas como até excessivas, em todas ou na maioria das questões foi percebida por 73,8 \% dos respondentes.
Considerando -se as informações/instruções fornecidas para a resolução das questões ( Questão 6 ), a proporção que considerou que eram até excessivas ou suficientes em todas ou na maioria das questões foi maior entre os estudantes de Universidades, \textbf{totalizando} 84,2 \%.
Os CEFET/IF foram a Organização Acadêmica em que houve a menor proporção : 79,9 \%.
No Gráfico 4.23, é possível observar que a diferença entre as Universidades e os CEFET/IF é estatisticamente significativa ( ver também Tabela II.12, no Anexo II ).
O percentual de estudantes que consideraram que as informações/instruções fornecidas para a resolução das questões eram até excessivas ou suficientes em todas ou na maioria das questões foi maior entre os estudantes de cursos de IES Públicas ( 84,1 \% ), não havendo diferença estatisticamente significativa da proporção registrada em IES Privadas ( 82,3 \% ) ( ver Gráfico 4.24 e Tabela II.12, no Anexo II ). 4.5 DIFICULDADE ENCONTRADA AO RESPONDER À PROVA \textbf{Indagados} sobre as dificuldades com as quais se depararam ao responder à prova ( Questão 7 ), 13,8 \% dos estudantes apontaram o Desconhecimento do conteúdo.
Para 42,8 \%, a Forma diferente de abordagem do conteúdo foi indicada como dificuldade.
Já a Falta de motivação para fazer a prova foi a dificuldade apontada por 20,3 \% dos respondentes.
Considerando -se todo o Brasil, 16,6 \% dos respondentes afirmaram que não tiveram qualquer tipo de dificuldade para responder à prova ( Tabelas II.13 e II.14, no Anexo II ).
Nos Gráficos 4.25 a 4.28, são apresentados os percentuais de estudantes que apontaram o Desconhecimento do conteúdo como “ dificuldade ao responder à prova ”.
Na análise por Grandes Regiões, o percentual de inscritos e presentes que apontaram o Desconhecimento do conteúdo como “ dificuldade ao responder à prova ” não superou 15,6 \%.
Os percentuais variaram de 12,1 \%, na região Sudeste, a 15,6 \%, na região Sul.
A Forma diferente de abordagem do conteúdo foi a escolha modal dos estudantes, com percentuais que variaram de 39,2 \% ( região Sudeste ) a 46,7 \% ( região Norte ).
O percentual de alunos que citaram a Falta de motivação como dificuldade variou de 17,5 \% ( região Sul ) a 25,0 \% ( região Centro-Oeste ).
Os que declararam não ter qualquer tipo de dificuldade para responder à prova variaram de 12,0 \%, na região Norte, a 19,2 \%, na região Sudeste.
Em relação aos quartos de desempenho, o Desconhecimento do conteúdo foi a opção escolhida por 16,6 \% dos estudantes do quarto inferior e por 11,1 \% do quarto superior, caracterizando uma tendência decrescente com o desempenho.
São observadas diferenças estatisticamente significativas entre o quarto inferior e os dois quartos superiores de desempenho, assim como entre o segundo quarto e o último quarto.
A alternativa modal para os alunos, quando agregados pelos quartos de desempenho, para a dificuldade encontrada, foi causada pela Forma diferente de abordagem do conteúdo : 38,6 \%, no último quarto, e 46,1 \% do terceiro quarto, escolheram essa alternativa.
Na análise por Organização Acadêmica, o percentual de inscritos e presentes que apontaram o Desconhecimento do conteúdo como “ dificuldade ao responder à prova ” não superou 17,3 \%.
Os percentuais variaram de 11,6 \%, nos CEFET/IF, a 17,3 \%, nos Centros Universitários.
A Forma diferente de abordagem do conteúdo foi a escolha modal dos estudantes, com percentuais que variaram de 42,5 \% ( Universidades ) a 45,8 \% ( CEFET/IF ).
O percentual de alunos que citaram a Falta de motivação como dificuldade variou de 11,6 \% ( Faculdades ) a 21,6 \% ( Universidades ).
Os que declararam não ter qualquer tipo de dificuldade para responder à prova variaram de 16,2 \%, nas Universidades e CEFET/IF, a 20,3 \%, nas Faculdades.
No Gráfico 4.27, é possível observar que há diferenças estatisticamente significativas dos Centros Universitários em relação a Universidades e a CEFET/IF ( ver também Tabela II.14, no Anexo II ).
Em relação às Categorias Administrativas, o Desconhecimento do conteúdo foi a opção escolhida por 16,3 \% dos estudantes nas IES Privadas e por 12,6 \% nas IES Públicas.
São observadas diferenças estatisticamente significativas entre essas duas categorias.
A alternativa modal entre os alunos, quando agregados pelas Categorias Administrativas, para a dificuldade encontrada, foi causada pela Forma diferente de abordagem do conteúdo : 43,4 \%, entre os estudantes de IES Privadas, e 42,4 \%, entre os de IES Públicas ( ver Gráfico 4.28 e Tabela II.14, no Anexo II ). 4.6 CONTEÚDOS DAS QUESTÕES OBJETIVAS DA PROVA Ao se analisarem os conteúdos das questões objetivas da prova ( Questão 8 ), verifica -se que um percentual muito pequeno dos estudantes avaliados, apenas 7,6 \%, afirmou que não estudou ainda a maioria desses conteúdos ( Gráficos 4.29 a Gráfico 4.32, e as Tabelas II.15 e II.16, no Anexo II ).
A maioria ( 64,7 \% ) afirmou ter estudado e aprendido muitos ou todos os conteúdos.
Na análise por Grande Região, a proporção de respondentes que escolheram a opção não estudou ainda a maioria desses conteúdos foi pequena.
Observa -se que, apenas na região Sul ( 11,6 \% ), a proporção foi maior do que a da média nacional ( 7,6 \% ).
Nas demais regiões Norte, Nordeste, Sudeste e Centro-Oeste, a situação foi \textbf{inversa}, com proporções mais baixas que a média, correspondendo a, respectivamente, 6,7 \%, 6,7 \%, 5,1 \% e 6,2 \%.
A diferença entre a região Sul e as demais é estatisticamente significativa, assim como entre as regiões Nordeste e Sudeste.
Em todas as regiões, a maioria dos presentes afirmou ter estudado e aprendido muitos ou todos os conteúdos, com proporções variando entre 57,4 \%, na região Sul, a 73,5 \%, na região Sudeste.
Considerando -se separadamente as opiniões de estudantes dos quatro quartos de desempenho, observa -se que, no quarto inferior, 14,2 \% ofereceram como resposta que não estudou ainda a maioria desses conteúdos, sendo 1,8 \% os do quarto superior com a mesma resposta.
A diferença entre os alunos que optaram por esse motivo de dificuldade em todos os quartos de desempenho é estatisticamente significativa.
Tendo -se em conta o quarto superior, 86,1 \% dos alunos afirmaram ter estudado e aprendido muitos ou todos os conteúdos.
No outro extremo, no primeiro quarto, 44,5 \% optaram pelas mesmas categorias.
As proporções são crescentes com o desempenho.
Na análise por Organização Acadêmica, a proporção de respondentes que escolheram a opção não estudou ainda a maioria desses conteúdos foi pequena.
Observa -se que as proporções nas Faculdades, Universidades e CEFET/IF, com 5,7 \%, 5,8 \% e 6,4 \%, respectivamente, foram menores do que a média nacional ( 7,6 \% ).
Nos Centros Universitários, a situação foi \textbf{inversa} : a proporção foi mais alta ( 14,1 \% ).
No Gráfico 4.31, é possível observar que as diferenças entre os Centros Universitários e os demais tipos de Organização Acadêmica são estatisticamente significativas.
Em todos os tipos de Organização Acadêmica, a maioria dos presentes afirmou ter estudado e aprendido muitos ou todos os conteúdos, com proporções variando entre 53,9 \%, nos Centros Universitários, a 68,9 \%, nos CEFET/IF ( ver também Tabela II.16, no Anexo II ).
Em relação às Categorias Administrativas, observa -se que, nas IES Públicas, 5,5 \% ofereceram como resposta que não estudou ainda a maioria desses conteúdos, sendo 11,6 \% nas IES Privadas com a mesma resposta.
A diferença entre os alunos que optaram por esse motivo de dificuldade nas duas categorias é estatisticamente significativa.
Considerando -se as IES Públicas, 68,7 \% dos alunos afirmaram ter estudado e aprendido muitos ou todos os conteúdos, sendo de 57,1 \% a proporção dos que escolheram essa resposta nas IES Privadas ( ver Gráfico 4.32 e Tabela II.16, no Anexo II ). 4.7 TEMPO GASTO PARA CONCLUIR A PROVA Ao responderem sobre o tempo de conclusão da prova ( Questão 9 ), perto de três quartos dos estudantes ( 73,4 \% ) afirmaram ter gasto Entre duas e quatro horas ( Gráfico 4.33 a Gráfico 4.36 e, no Anexo II, as Tabelas II.17 e II.18 ).
Considerando -se as cinco Grandes Regiões brasileiras, os percentuais dos que utilizaram Entre duas e quatro horas para finalizar a prova nas regiões Norte ( 71,3 \% ), Sudeste ( 72,6 \% ) e Centro-Oeste ( 70,8 \% ) foram inferiores ao percentual nacional.
Nas demais Grandes Regiões, os percentuais de alunos que \textbf{dispensaram} Entre duas e quatro horas para concluir a prova ficaram muito próximos : 74,0 \%, na região Nordeste, e 74,3 \%, na região Sul, como mostrado no Gráfico 4.33.
Não há diferenças estatisticamente significativas entre as regiões.
Uma vez analisadas as alternativas escolhidas pelos estudantes que se situam nos diferentes quartos de desempenho, observa -se uma tendência crescente : uma proporção maior de participantes no quarto superior declarou ter gasto Entre duas e quatro horas para concluir a prova quando comparada com a proporção de participantes nos quartos inferiores.
Há diferenças estatisticamente significativas entre o quarto inferior e os demais quartos de desempenho, bem como entre o segundo e o quarto superior, e também entre o terceiro e o quarto quartos de desempenho.
Na análise por Organização Acadêmica, o percentual dos que utilizaram Entre duas e quatro horas para finalizar a prova nas Universidades ( 73,9 \% ) foi ligeiramente superior ao percentual nacional ( 73,4 \% ).
Nas Faculdades, nos Centros Universitários e nos CEFET/IF, os percentuais foram quase os mesmos, respectivamente, 72,2 \%, 72,0 \% e 72,0 \%.
Como mostrado no Gráfico 4.35, não há diferença estatisticamente significativa entre as proporções dos quatro tipos de Organização Acadêmica.
Em relação às Categorias Administrativas, observa -se que, nas IES Públicas, 73,4 \% declararam ter gasto Entre duas e quatro horas para concluir a prova, sendo 73,2 \% nas IES Privadas com a mesma resposta.
A diferença entre os alunos que optaram por esse motivo de dificuldade nas duas categorias não é estatisticamente significativa ( ver Gráfico 4.36 e Tabela II.18, no Anexo II ).

% matched lemmas: arredondamento, crescer, disjunto, dispensar, estimador, estudantes, indagar, interseção, inverso, numérico, ordenar, parâmetro, percentagem, questionário, recorte, reprodução, subconjunto, sumário, totalizar, vidar
\end{textsample}
